% generated by qpp/qpp2latex.py from qpp/examples/means.qpp; do not edit

Elementary comparisons of masses.

\section{Means and squares}\label{sec:means}

\begin{fact}[{power means, \cite{HornJohnson2013}}]\label{fact:powermean}
The power mean inequality at the exponent $E + 1$: H\"older's inequality against the vector of ones, with the
source's constant $J^E$.
\begin{stmt}
\sit{given}{J,E\in \N_{\ge1},\ u_1,\dots,u_J\in \Q_{>0}}
\sit{(1)}{(\textstyle\sum_{j=1}^J u_j)^{E+1} \le J^E\cdot (\textstyle\sum_{j=1}^J u_j^{E+1})}
\end{stmt}
\end{fact}

\begin{lemma}[{the mean of two masses}]\label{lem:mean}
The mean of two masses lies strictly between them.
\begin{stmt}
\sit{given}{u,v\in \Q_{>0},\ u < v}
\sit{(1)}{u < (u+v)/2}
\sit{(2)}{(u+v)/2 < v}
\end{stmt}
\end{lemma}

\begin{proof}
Half the gap on each side.
\begin{plines}
\pline{(1)}{The goal follows from \(1/2\) times \((\mathrm{H}_{1})\).}
\pline{(2)}{The goal follows from \(1/2\) times \((\mathrm{H}_{1})\).}
\end{plines}
\end{proof}

\begin{lemma}[{the square of a sum}]\label{lem:square}
The square of a sum, with its cross term named.
\begin{stmt}
\sit{given}{u,v\in \Q_{>0},\ c:=2 u v}
\sit{(1)}{(u+v)^2 = u^2+v^2+c}
\sit{(2)}{c < (u+v)^2}
\end{stmt}
\end{lemma}

\begin{proof}
The cross term is a name: the identity uses its definition, never its expansion.
\begin{plines}
\pline{(1)}{The goal follows from \((\mathrm{H}_{c})\).}
\pline{(2)}{The goal follows from \((\mathrm{H}_{c})\).}
\end{plines}
\end{proof}

\begin{lemma}[{a sum of squares}]\label{lem:sumsq}
The square of a sum of masses is at most their number times the sum of their squares.
\begin{stmt}
\sit{given}{J\in \N_{\ge1},\ u_1,\dots,u_J\in \Q_{>0}}
\sit{(1)}{(\textstyle\sum_{j=1}^J u_j)^2 \le J\cdot (\textstyle\sum_{j=1}^J u_j^2)}
\sit{(2)}{\text{equality exactly when the}\ u_j\ \text{are all equal}}
\end{stmt}
\end{lemma}

\begin{proof}
The power mean inequality at the exponent two.
\begin{plines}
\pline{(1)}{The goal follows by \cref{fact:powermean}(1).}
\end{plines}
\end{proof}
